\documentclass[12pt]{article}

\usepackage[a4paper, total={6in, 8in}]{geometry}
\usepackage{textcomp}

\begin{document}
\setlength{\parindent}{0pt}

\def \t {\textrightarrow}
\def \v {\vspace{1em}}
\def \bi {\begin{itemize}}
\def \ei {\end{itemize}}
\def \s[#1] {\section*{#1}}
\def \ss[#1] {\subsection*{#1}}
\def \sss[#1] {\subsubsection*{#1}}

\s[Nevicata - Carducci]
Descrive l'immagine di questo vetro che costituisce un limite \t e il limite del vetro è quello di vedere una realtà, ma non poterla abbracciare \\
Si trovano aspetti umani e autobiografici, che assumono qui uno spessore decadente \\
Atmosfera ovattata di un mattino di nebbia, e la neve scende lenta e fioca (metafora perchè fa i fiocchi, e analogia) \\
Importante è anastrofe tra soggetto (neve) e fioca \\
Cinereo \\
Dopo i due punti \\
Verso due \t iperbato, e mira a creare atmosfera plumbea e ovattata \\
Il grido d'erbaiola è iperbato \t richiamo leopardiano \t non e non = anafora \\
Il borgo, le figure indefinite della figura leopardiana \t Petrarca silenzio della poesia, il viaggio dell io si passa da Leopardi \t in Pascoli viaggio interiore è parziale \\
La torre rimanda al passero solitario \\
Roche = canto ingozzato, lo scandirsi del suono delle ore \t che evoca un rantolo quasi, che è di per se gia un segno di morte \\
Vetri appannati, e gli uccelli volteggiano \t contro cui gli uccelli sbattono \t sono animali di solitudine \t richiamo di un altro mondo dietro un velo, che sono i vetri appannati \t un limite \\
Spezzatura sonora con gli ultimi due versi \t indomito cuore anche questa è un'evocazione leopardiana (in A se stesso, leopardi si rivolge direttamente al suo cuore) \\
Indomito significa che non si riesce a controllare \\
Calmati \t imperativo \t discorso diretto libero \\
Ultimo verso è un chiasmo, silenzio e ombra \t connotati dell'oltretomba \t evocazione di versi danteschi "sovra lor vanita che par persona" (nel mondo degli ignavi)

\v

La lirica presenta una trama narrativa esile \t struttura pulita, e una abbondanza di immagini evocative \\
Questo paesaggio da idea di cio che i famosi fiamminghi (i 3, rivedere) \t mettono a fuoco la letteratura decadente, in italia oltrepassata attraverso i "Crepuscolari" ???
\end{document}
